\section{Related Work}
\label{sec:related}

\paragraph{Memory poisoning, and where its defences sit.}
\citet{mpbench2026} systematize memory-poisoning attacks on LLM agents
and report the split that motivates this work: every off-the-shelf
detector they test loses between $18$ and $42$ points of true-positive
rate moving from strong-signal payloads to weak-signal ones---payloads
written to read like ordinary operational facts, carrying, in their
reading, ``no syntactic anomaly''. \citet{memsecbench2026} track poisoning through a lifecycle
and identify selective repair as the open stage: removing poison without
destroying benign memory alongside it. We adopt their
Write--Execute--Forget checkpoint structure as a protocol
(\S\ref{sec:wef}), reimplemented rather than ported, and make no
comparison to their numbers.

What both share with the wider defence literature is \emph{where} the
defence sits: at write, screening text before it is stored, or at
retrieval, screening entries before they are injected. Both read content,
and content is exactly what a weak-signal payload is constructed to
survive. Our threat model is orthogonal to that axis rather than a better
position on it: we do not ask whether the defence can recognize the
payload, but what it does when the evidence it decides on has itself been
tampered with.

\paragraph{Attacking the curator: named, and left unmeasured.}
We are not the first to describe this threat, and we are explicit about
that because the distinction between naming it and measuring it is the
whole of our contribution here. \citet{lin2026selfevolving} catalogue
the attack surface of self-evolving agent systems and state the model
directly: an adversary who can ``shift embedding proximity or corrupt
the performance signal gains indirect control over the long-term
composition of the agent's cognitive state without ever writing to the
memory store directly,'' with the consequence that ``biased selection
can systematically retain harmful memories while discarding
safety-critical ones.'' Their case study CS-I6 names the same mechanism
as \emph{selective amplification}. \citet{lin2026memsurvey} organize
memory security around a lifecycle whose final phase is
\emph{Forget \& Rollback}, and place denial-style consequences---silent
eviction under retention pressure, denial of retrieval---in its
availability column.

Both also record that the threat is unmeasured. \citet{lin2026selfevolving}
write that ``direct attacks on memory consolidation policies remain
relatively underexplored''; \citet{lin2026memsurveyv1} report that
availability threats on agent memory are ``architecturally plausible but
empirically unstudied''. \pointer{We cite that survey's first version
advisedly, and \S\ref{sec:citationnote} says why.}
\inplace{related-survey-versions}
That is the gap this paper closes. Our contribution is not the
observation that a curator can be attacked through its evidence; it is
the operationalization---a resource-bounded adversary, six curation
mechanisms, five attack budgets, thirty seeds, per-seed intervals, exact
paired permutation tests, and a published failure boundary for the
mechanism we defend with---together with a real-task leg that reports
which half of that result transfers and which does not
(\S\ref{sec:swebench-attack}).

\paragraph{The nearest neighbour, and the cell it leaves open.}
\citet{wu2026forget} curate on-device agent memory under a hard byte and
energy budget, which is a genuine conserved constraint --- but the eviction
decision is a score, ``value minus harm, per byte'', a learned-utility ruler
applied to each entry. That is the axis this paper turns on: a budget is not
enough if a valuation still decides what the budget spends on. Our entries
are not scored for value; they hold a balance that only a measured
environment outcome moves, so the thing evicted is the thing that failed to
earn, not the thing a ruler judged low. Conserved-and-scored and
judged-but-unbudgeted are the two cells adjacent to ours; the unclaimed one
is conserved \emph{and} settled by measurement \emph{and} judge-free at once.


\pointer{Every other neighbour sits in a different cell, and
\S\ref{sec:neighbours} places each: memory-as-a-model
\citep{quek2026memo}; environment-mediated selection over behaviours
\citep{dodgson2026survival}; hand-designed salience, paging and verbal
accumulation \citep{park2023generative,packer2023memgpt,
shinn2023reflexion,zhao2024expel,wang2023voyager}; selection over agents
\citep{zhang2025dgm}; retrieval-policy bandits \citep{xu2026ael};
LLM-curated banks and their decay \citep{zhang2026faulty}; the
continual-learning testbed and the objection bounding it
\citep{jimenez2024swebench,joshi2025swebenchcl,agentcl2026};
representation as the orthogonal axis
\citep{reasoningbank2026,cl2memory2026}; attacks that never write
\citep{chen2024agentpoison,zou2026poisononce,srivastava2025memorygraft,xu2026steering,dong2025minja};
admission and earned-versus-spent budgets
\citep{admission2026,ember2026,tmem2026,xu2025amem,mem0}; and learned
credit assignment \citep{memoryr2,treecredit2026,attrimem2026}.}
\inplace{related-neighbours}
